\documentclass[11pt]{article}

% ---------------------------------------------------------------------------
%  L2 — The mod-3 arithmetic of the polyplet height triangle.
%  Category L (entirely machine-written); see paper/README.md and
%  docs/publication-split.md.
%
%  NOVELTY STATUS: UNCHECKED.  The N1-N6 sweeps covered the diagonal law,
%  haruspicy, the staircase squeeze, HV-convex polyplets and the lambda
%  upper bound.  They never touched the ternary spine, the digit product or
%  the SNF count.  docs/publication-split.md §5 requires a sweep in the
%  N1-N6 style BEFORE this is written -- it has not happened, and the paper
%  says so in three places rather than one.
%
%  Source material, all in this repository:
%    results/ternary-spine.md      T1-T5, the ladder, the row reading
%    results/triangle-snf.md       the SNF thread this closes
%    results/defect-gas.md         the master equation and the weights
%    docs/proofs/diagonal-law.md   the law this all rests on
%    experiments/ternary_spine.py  15/15 verifier
% ---------------------------------------------------------------------------

% ---------------------------------------------------------------------------
%  paper/shared/preamble.tex — packages and macros common to every manuscript
%  in this directory.  Factored out of paper/polyplets-report.tex and
%  paper/technical-report.tex, which between them had two copies of the same
%  twenty lines.
%
%  technical-report.tex does NOT \input this file and is not to be edited to do
%  so: it is jasonp's prose and is read-only to the machine (paper/README.md,
%  "Who may edit what").  The duplication there is deliberate.
%
%  Assumes \documentclass[11pt]{article} has already been issued.
% ---------------------------------------------------------------------------
\usepackage[margin=1in]{geometry}
\usepackage{amsmath,amssymb}
\usepackage{amsthm}
\usepackage{booktabs}
\usepackage{graphicx}
\usepackage{numprint}
\usepackage{microtype}
\usepackage[table]{xcolor}
\usepackage[hidelinks]{hyperref}
\usepackage{flafter}   % a float never appears before the text that introduces it
\usepackage{float}     % [H] pins tables/figures exactly where written
\usepackage[font=small,labelfont=bf,labelsep=period]{caption}
\setlength{\intextsep}{16pt plus 3pt minus 3pt}

\npthousandsep{,}
\npfourdigitsep
\newcommand{\num}[1]{\numprint{#1}}
\newcommand{\g}{\cellcolor{gray!15}}

% Theorem environments.  One counter shared across kinds, so that a reference
% like "Theorem 4" is unambiguous in a paper that also has lemmas and
% propositions.
\theoremstyle{plain}
\newtheorem{theorem}{Theorem}
\newtheorem{proposition}[theorem]{Proposition}
\newtheorem{lemma}[theorem]{Lemma}
\newtheorem{corollary}[theorem]{Corollary}
\theoremstyle{definition}
\newtheorem{definition}[theorem]{Definition}
\newtheorem{example}[theorem]{Example}
\theoremstyle{remark}
\newtheorem{remark}[theorem]{Remark}
\newtheorem{conjecture}[theorem]{Conjecture}
\newtheorem{openproblem}[theorem]{Open Problem}

% Repo-path citations.  Every one of these must resolve in the tree that ships
% with the paper; `make gate-citations` checks the markdown, and
% paper/verify_claims.py checks the manuscripts.
\newcommand{\repo}[1]{\texttt{#1}}

% OEIS links.
\newcommand{\oeis}[1]{\href{https://oeis.org/#1}{#1}}

% Growth constants, so a rename is one edit.
\newcommand{\lam}{\lambda}
\newcommand{\muH}{\mu_H}

% ---------------------------------------------------------------------------
%  paper/shared/disclosure.tex — the two authorship disclosure blocks.
%
%  docs/publication-split.md §1 fixes the rule these implement:
%
%    "Under no circumstances will there be any ambiguity about whether a paper
%     is wholly my work (none), merely my prose (some) or entirely LLM-authored
%     (some)."   — jasonp, 2026-08-07
%
%  There are exactly two blocks and they live in exactly one file, so that a
%  paper cannot quietly soften its own disclosure.  The fixed sentences take no
%  arguments.  The ONE thing a paper supplies is the per-result verification
%  ledger, because that is the part that differs honestly between papers — and
%  §1 requires it to be per result, not in aggregate.
%
%  Do not add a \Pdisclosure or \Ldisclosure variant with an optional softening
%  argument.  If a paper does not fit one of these two blocks, the paper is
%  wrong, not the block.
% ---------------------------------------------------------------------------

\newsavebox{\disclosurebox}

\newenvironment{disclosureframe}
  {\begin{center}\begin{lrbox}{\disclosurebox}\begin{minipage}{0.93\textwidth}\small}
  {\end{minipage}\end{lrbox}\fbox{\usebox{\disclosurebox}}\end{center}}

% --- P: jasonp's prose -------------------------------------------------------
% Byline: Jason Parker, alone.  Argument: the per-result ledger of what the
% machine did.
\newcommand{\Pdisclosure}[1]{%
\begin{disclosureframe}
\textbf{Authorship.}\enspace Every sentence of this paper was written by its
named author. He has read every proof it states and believes each one, and he
is responsible for the correctness of what it claims.

\smallskip
\textbf{Machine contribution.}\enspace #1
\end{disclosureframe}}

% --- L: entirely machine-written --------------------------------------------
% Argument: the per-result verification ledger — for each result, what warrants
% it (Lean kernel, exact certificate, or labelled measurement) and how much of
% it a human has checked.
\newcommand{\Ldisclosure}[1]{%
\begin{disclosureframe}
\textbf{Authorship.}\enspace This document was written end to end by a large
language model. That includes its mathematics, its prose, its tables and its
\LaTeX{}. No sentence of it was written by a human. The mathematics it reports
was likewise derived by machine.

\smallskip
\textbf{Human role.}\enspace The person who directed this work chose what to
compute, chose what to publish, and is responsible for the decision to release
this document. Except where the ledger below says otherwise, that person has
not personally checked the arguments here, and this disclosure is not to be
read as an endorsement of them.

\smallskip
\textbf{What warrants each result.}\enspace #1
\end{disclosureframe}}

% --- Draft banner ------------------------------------------------------------
% For an L-paper whose verification ledger is not yet settled.  Loud on purpose:
% an L-paper without a truthful ledger must not be mistaken for a finished one.
%
% \firstdraftbanner is the standard wording, which five papers previously
% carried character-for-character in their own preambles.  It lives here for the
% same reason the disclosure blocks do: identical text in seven places drifts,
% and a paper must not be able to soften it by editing its own copy.  The
% optional argument is for a gate specific to one paper, appended verbatim; a
% paper needing different wording (L6, compute-gated) calls \draftbanner direct.
\newcommand{\firstdraftbanner}[1][]{%
\draftbanner{This is a first draft. The verification ledger below records what
warrants each result \emph{mechanically}; the column that records what a human
has read is empty, deliberately, because nobody has read it yet. Do not
circulate until that column says something true.#1}}

\newcommand{\draftbanner}[1]{%
\begin{center}
\fcolorbox{red!60!black}{yellow!15}{%
  \begin{minipage}{0.93\textwidth}\small
  \textbf{DRAFT — NOT FOR CIRCULATION.}\enspace #1
  \end{minipage}}
\end{center}}


\newcommand{\Z}{\mathbb{Z}}
\newcommand{\Q}{\mathbb{Q}}
\newcommand{\F}{\mathbb{F}}
\newcommand{\vthree}{v_3}

\title{\textbf{The mod-3 arithmetic of the polyplet height triangle}\\[4pt]
\large One cubic, a base-3 digit product, and a Smith normal form}
\author{[byline: jasonp's call --- \repo{docs/publication-split.md} \S1]}
\date{August 2026}

\begin{document}
\maketitle

\firstdraftbanner[ A second gate applied here and is now cleared. The paper was
drafted before its literature sweep, which \repo{docs/publication-split.md} \S5
forbids; the sweep ran on 2026-08-18 and found no collision
(\S\ref{sec:novelty}).]

\Ldisclosure{%
\emph{Theorem~\ref{thm:snf-3power} (all invariant factors are $3$-powers)} ---
unconditional, two lines, proved below.
\emph{Propositions~\ref{prop:polya} and \ref{prop:GH}} --- unconditional given
the diagonal law.
\emph{Theorems~\ref{thm:digit}, \ref{thm:oddspine}, \ref{thm:evenspine},
\ref{thm:activation}, \ref{thm:snf}} --- conditional on the diagonal law
(itself a theorem, companion paper L1) together with the three congruences of
\S\ref{sec:ladder}, each derived symbolically from the cluster master equation. All
five are additionally verified against an independently enumerated triangle:
$342$ in-regime cells for the digit product, all $k \le 17$ for the spine
theorems, and every $N \le 36$ for the Smith normal form.
\emph{Theorem~\ref{thm:deficit2} (the deficit-$2$ residue)} --- proved by
Lagrange--B\"urmann over $\Z/27$ followed by an exact polynomial division on
the mod-$27$ master curve; the division is a finite exact computation.
\emph{Theorem~\ref{thm:d3} (the deficit-$3$ residue)} --- proved on the frame
stated with it: (A1) and (A4) are proved and independently cross-checked, while
(A2) and (A3) are symbolic derivations verified against banked data to $k \le 17$
rather than proofs. Everything above the frame is exact finite algebra.
\emph{\S\ref{sec:sleeve} (sleeve zeros)} --- labelled measurement, complete for
$n \le 40$, with no closed form claimed \emph{except} at deficit $3$, where
Theorem~\ref{thm:d3} predicts them for all $k$.
\emph{Novelty} --- pass run 2026-08-18, no collision found on the spine cubic,
the digit product or the Smith normal form count; the two bodies of machinery a
reader will think of are named in \S\ref{sec:related}, with why neither applies.
\emph{Human verification:} none, as of this draft.}

\begin{abstract}
\noindent
Let $T(n,H)$ be the number of fixed king-lattice animals with $n$ cells whose
bounding box has height exactly $H$, and let $T_N = [T(i,j)]_{i,j \le N}$.
Because $T_N$ is lower triangular with determinant $3^{N(N-1)/2}$, its Smith
normal form over $\Z$ is $\mathrm{diag}(3^{e_1},\dots,3^{e_N})$; the number of
nontrivial invariant factors is $\lceil (N-1)/3\rceil$. We derive it from the
arithmetic of the triangle modulo $3$, assuming the diagonal law and the product
form of companion paper L1 together with three congruences derived from the
cluster master equation.

Modulo $3$ the diagonal generating series of the triangle is the unique
$W \in \F_3[[t]]$ with $W(0)=1$ satisfying
\[
  W^3 \;=\; W^2 + t .
\]
Writing $n = \sum_i n_i 3^i$ in base $3$, the diagonal polynomials satisfy
\[
  P_k(n) \;\equiv\; [y^k]\prod_i W\!\left(y^{3^i}\right)^{n_i} \pmod 3 ,
\]
so the triangle modulo $3$ is $3$-automatic with an explicit automaton.

The cubic admits the parametrisation
\[
  t \;=\; \frac{x}{(1-x)^3}, \qquad W \;=\; \frac{1}{1-x} ,
\]
in which the Lagrange--B\"urmann correction term vanishes identically in
characteristic $3$, so the two boundary cases can be evaluated in closed form:
column $H$ vanishes modulo $3$ in every row below $\lfloor 3H/2\rfloor$ and is a
unit in that row. Since $H \mapsto \lfloor 3H/2\rfloor$ is strictly increasing,
the surviving columns of $T_N$ modulo $3$ stand in echelon position, and
counting the vanishing ones gives the invariant-factor count.

Read along rows, the last nonzero residue of row $n$ is $1$, $1$, $2$ according
to $n \bmod 3$, the third case being the deficit-$2$ cell, which we treat
separately by Lagrange--B\"urmann over $\Z/27$. We give a complete census for
$n \le 40$ of the cells whose $3$-adic valuation exceeds the deficit the
diagonal law forces on them; a closed formula for those exceptions is open. The
cubic lifts: $H^3 - H^2 \equiv 25t \pmod{27}$ with $25$ the single-defect weight
of L1's cluster decomposition; $W$ returns as the level-$3$ correction.
\end{abstract}

\tableofcontents

\section{Introduction}

\subsection{The observation}

The height triangle of the king lattice has entry $T(n,H)$, the number of fixed
polyplets with $n$ cells whose bounding box has height exactly $H$. Its
leading $N \times N$ block $T_N$ is lower triangular, since $T(n,H) = 0$ for
$H > n$, and every invariant factor of its Smith normal form over $\Z$ is a
power of $3$.

The number of \emph{nontrivial} invariant factors, those with exponent $> 0$,
does not follow from the determinant, and is $\lceil (N-1)/3\rceil$:

\begin{table}[H]
\centering\small
\begin{tabular}{l r r r r r r r r r r r}
\toprule
$N$ & 4 & 5 & 6 & 7 & 8 & 9 & 10 & 11 & 12 & 13 & 14 \\
\midrule
nontrivial factors & 1 & 2 & 2 & 2 & 3 & 3 & 3 & 4 & 4 & 4 & 5 \\
\bottomrule
\end{tabular}
\caption{One new nontrivial invariant factor per three columns.}
\label{tab:snfcount}
\end{table}

\subsection{Novelty status}
\label{sec:novelty}

\textbf{Swept, no collision.} The project's first six literature sweeps covered
the diagonal law, haruspicy and non-D-finiteness, the staircase squeeze,
HV-convex polyplets and the growth-constant upper bound, and none of them
touched this paper's objects. A seventh ran on 2026-08-18 and is recorded in
\repo{docs/priority-passes-2026-08-18.md}: congruences modulo powers of $3$ for
combinatorial sequences, the Krattenthaler--M\"uller method for mod-$3^k$
behaviour of recursive sequences, Rowland--Yassawi automatic congruences for
diagonals of rational functions, Smith normal forms of combinatorial triangles,
and base-$3$ digit-product formulas. It found no collision on the spine cubic,
the digit product or the Smith normal form count. What it did find is one
adjacent method, Rowland--Yassawi, treated in \S\ref{sec:related}. The nearest
neighbours otherwise are Christol's theorem, which makes ``algebraic over
$\F_3$'' and ``$3$-automatic'' the same statement, and the theory of
Artin--Schreier curves, to which the cubic below belongs.

\section{Setup}
\label{sec:setup}

We take as given the \emph{diagonal law}, which is a theorem: for every $k$
there is a polynomial $P_k$ of degree exactly $k$ with
\begin{equation}
  T(n,\,n-k) \;=\; P_k(n)\cdot 3^{\,n-1-3k}
  \qquad\text{for every } n \ge 2k+1,
  \label{eq:law}
\end{equation}
and the \emph{product form}, also a theorem: there are $G, H \in \Q[[y]]$ with
$G(0) = H(0) = 1$ and
\begin{equation}
  \sum_{k \ge 0} P_k(n)\,y^k \;=\; G(y)\,H(y)^n
  \qquad\text{for every } n \ge 2k+1 \text{ at order } y^k .
  \label{eq:grand}
\end{equation}
Both are proved in companion paper L1, \emph{A diagonal law for row-local
lattices}. Equation \eqref{eq:grand} is what makes this paper possible, since it
replaces infinitely many polynomials by two power series. The cubic of
\S\ref{sec:cubic} comes from the same source: L1's \emph{universal spine
theorem} gives that one curve for every prime $p$ dividing the base, provided
the lattice's single-defect weight is nonzero modulo $p$ and $t$ is rescaled by
it. Here $p = 3$ and the weight is $25 \equiv 1 \pmod 3$, so no rescaling is
visible.

\begin{proposition}[P\'olya integrality]
\label{prop:polya}
Each $P_k$ is integer-valued, hence $3$-integral: $\vthree(P_k(n)) \ge 0$ for
every integer $n$.
\end{proposition}

\begin{proof}
The window $[2k+1,\,3k+1]$ contains exactly $k+1$ consecutive integers, and on
it $P_k(n) = T(n,n-k)\cdot 3^{\,3k+1-n}$ is a nonnegative count times a
nonnegative power of $3$, hence an integer. A polynomial of degree $k$ that is
integral at $k+1$ consecutive integers is integer-valued
everywhere~\cite{polya1915}.
\end{proof}

The leading coefficient of $P_k$ is $25^k/k!$, so $P_k$ takes integer values
without having integer coefficients, as $\binom{n}{k}$ does; the window in the
proof is the $k+1$ consecutive integers such a polynomial needs.

\begin{proposition}
\label{prop:GH}
$G = \sum_k P_k(0)y^k$ and $H = \bigl(\sum_k P_k(1)y^k\bigr)/G$ are integer power
series with constant term $1$.
\end{proposition}

\begin{proof}
Evaluate \eqref{eq:grand} at $n = 0$ and $n = 1$. Both left sides are integral
by Proposition~\ref{prop:polya}; $G$ has constant term $P_0(0) = 1$, so it is
invertible in $\Z[[y]]$ by the geometric series.
\end{proof}

Proposition~\ref{prop:GH} computes $H$ from the diagonal polynomials, and with
$P_k$ known through $k = 17$ it begins $1, 25, 208, 1483, 20688, 130208, \dots$
It is the per-row transfer series of L1's cluster decomposition, and its first
coefficient $25$ counts the single-defect species of one row, which is L1's pair
weight. Whether $H$ is a known sequence is for the missing novelty sweep
to settle; it has not been looked up in the OEIS~\cite{oeis}.

\section{Three coefficient identities}
\label{sec:ladder}

Everything below uses three congruences, each derived from L1's cluster master
equation
\[
  H \;=\; 1 + \sum_c \widehat W_c\, u^{k_c} H^{-(k_c+\ell_c)} ,
\]
the sum running over cluster types $c$ of surplus $k_c$ and length $\ell_c$.

\begin{description}
\item[$(\star a)$] $G \equiv 1 \pmod 9$, from a residue formula for the boundary
  clusters, a valuation lemma on boundary weights, and one identity of the
  mod-$9$ cubic.
\item[$(\star b)$] $H \equiv W \pmod 3$. Every cluster has surplus at least its
  length, so the valuation lemma leaves only the two-cell row surviving modulo
  $3$: the cubic counts chains of those rows and nothing else. The derivation
  holds at all orders.
\item[$(\star c)$] $(H^3 - H(t^3))/3 \equiv t^2 + tW \pmod 3$, from the mod-$9$
  cubic alone. Uniqueness of its solution gives
  $H(t^3) \equiv H^2 + 25t - 3t^2 - 3tW \pmod 9$, and a three-line cancellation
  over $\F_3$ finishes.
\end{description}

$(\star b)$ reproduces $342$ independently enumerated triangle cells through the
digit product of \S\ref{sec:cubic}, and $(\star c)$ holds to order $t^{300}$
against the mod-$27$ solution of \S\ref{sec:lift}.

\section{The spine cubic and the digit product}
\label{sec:cubic}

\begin{definition}
$W \in \F_3[[t]]$ is the unique root with $W(0) = 1$ of $W^3 = W^2 + t$.
\end{definition}
Written $W^2(W-1) = t$, the cubic determines the coefficients of
$W = 1 + \dots$ one after another, which gives existence and uniqueness; $W$ is
algebraic of degree $3$ over $\F_3(t)$. Every computation below goes through its
rational parametrisation
\begin{equation}
  t \;=\; \frac{x}{(1-x)^3},
  \qquad
  W \;=\; \frac{1}{1-x} ,
  \label{eq:param}
\end{equation}
checked by $W^3 - W^2 = (1-x)^{-3}\bigl(1 - (1-x)\bigr) = x(1-x)^{-3} = t$.

\begin{theorem}[digit product]
\label{thm:digit}
Write $n = \sum_i n_i 3^i$ in base $3$. Then
\[
  P_k(n) \bmod 3 \;=\; [y^k]\;\prod_i W\!\left(y^{3^i}\right)^{n_i}.
\]
\end{theorem}

\begin{proof}
By $(\star a)$ and $(\star b)$, the product form \eqref{eq:grand} reduces modulo
$3$ to $\sum_k P_k(n)y^k \equiv W^n$. Frobenius over $\F_3$ gives
$W(y)^{3^i} = W(y^{3^i})$, so $W^n = \prod_i \bigl(W^{3^i}\bigr)^{n_i} =
\prod_i W(y^{3^i})^{n_i}$.
\end{proof}

Since $W$ is algebraic over $\F_3(t)$, Christol's
theorem~\cite{christol1980} makes the array $P_k(n) \bmod 3$
$3$-automatic~\cite{allouche2003}, and Theorem~\ref{thm:digit} writes that
automaton out in two dimensions.

\section{The spine}

Call the cell $T\bigl(n, \lfloor 2n/3 \rfloor\bigr)$ at
$n = \lfloor 3H/2\rfloor$ the \emph{spine} of column $H$; it is the first entry
of that column not divisible by $3$. Odd and even $H$ need separate treatment,
because the exponent $n-1-3k$ of \eqref{eq:law} is $0$ at the spine in the first
case and $-1$ in the second.

Both proofs use Lagrange--B\"urmann in the following form~\cite[Ch.~5]{stanley2012}:
if $x = t\varphi(x)$ then
\begin{equation}
  [t^k]\,F(x(t)) \;=\; [x^k]\;F(x)\,\varphi(x)^k
  \left(1 - \frac{x\varphi'(x)}{\varphi(x)}\right),
  \label{eq:LB}
\end{equation}
valid over any commutative ring. With the parametrisation \eqref{eq:param},
$\varphi(x) = (1-x)^3$, so
\[
  \varphi'(x) \;=\; -3(1-x)^2 \;=\; 0 \quad\text{in characteristic } 3 .
\]
The correction factor in \eqref{eq:LB} is therefore identically $1$.

\begin{theorem}[odd spine]
\label{thm:oddspine}
$T(3k+1,\,2k+1) \equiv 1 \pmod 3$ for every $k \ge 0$.
\end{theorem}

\begin{proof}
At $n = 3k+1$ the exponent in \eqref{eq:law} is $n-1-3k = 0$, so
$T(3k+1,2k+1) = P_k(3k+1) \equiv [t^k]W^{3k+1} \pmod 3$ by
Theorem~\ref{thm:digit}. Apply \eqref{eq:LB} with $F = (1-x)^{-(3k+1)}$:
\[
  [t^k]W^{3k+1} \;=\; [x^k]\,(1-x)^{-(3k+1)}(1-x)^{3k}
  \;=\; [x^k](1-x)^{-1} \;=\; 1 .
  \qedhere
\]
\end{proof}

\begin{theorem}[even spine]
\label{thm:evenspine}
$P_k(3k) \equiv 3 \pmod 9$ for every $k \ge 1$; equivalently
$T(3k,2k) = P_k(3k)/3 \equiv 1 \pmod 3$.
\end{theorem}

\begin{proof}
Let $S := (H^3 - H(t^3))/3$, integral because $H^3 \equiv H(t^3) \pmod 3$ is the
classical Frobenius congruence for integer series.

Suppose first $3 \nmid k$. Then
$H^{3k} \equiv H(t^3)^k + 3k\,S\,H(t^3)^{k-1} \pmod 9$. The first term involves
only exponents divisible by $3$, so $[t^k]$ annihilates it. Using $(\star a)$,
$(\star c)$ and $W(t^3) = W^2 + t = W^3$,
\[
  P_k(3k) \;\equiv\; 3k\,[t^k]\bigl(t^2 + tW\bigr)W^{3k-3} \pmod 9 .
\]
By \eqref{eq:LB} again, whose correction term is $1$, together with
\eqref{eq:param},
\[
  [t^k]\bigl(t^2+tW\bigr)W^{3k-3}
  \;=\; [x^k]\left(\frac{x^2}{(1-x)^3} + \frac{x}{1-x}\right)
  \;=\; \binom{k}{2} + 1 ,
\]
and for $3 \nmid k$ one has
$k\bigl(\binom{k}{2}+1\bigr) \equiv k\cdot k \equiv 1 \pmod 3$, so
$P_k(3k) \equiv 3 \pmod 9$.

If $3 \mid k$, write $k = 3m$. The middle term now carries a factor
$3 \cdot 3m \equiv 0 \pmod 9$, leaving the self-similarity
\[
  P_{3m}(9m) \;\equiv\; [t^{3m}]H(t^3)^{3m}
  \;=\; [\tau^{m}]H(\tau)^{3m} \;\equiv\; P_m(3m) \pmod 9 ,
\]
and induction on $\vthree(k)$ reduces to the previous case.
\end{proof}

We have verified the self-similarity for $m \le 5$, and the congruence
$P_k(3k)\equiv3\pmod 9$ for all $k \le 17$, against the enumerated triangle.

\begin{theorem}[the activation boundary]
\label{thm:activation}
Column $H$ of the triangle satisfies $T(n,H) \equiv 0 \pmod 3$ for every
$n < \lfloor 3H/2\rfloor$, and $T\bigl(\lfloor 3H/2\rfloor, H\bigr)
\not\equiv 0 \pmod 3$.
\end{theorem}

\begin{proof}
With $k = n - H$, \eqref{eq:law} gives
$\vthree(T(n,H)) = (3H-2n-1) + \vthree(P_{n-H}(n))$.

\emph{Below the boundary.} $n < \lfloor 3H/2\rfloor$ makes the explicit exponent
$3H-2n-1 \ge 1$, and $\vthree(P_k(n)) \ge 0$ by Proposition~\ref{prop:polya}. So
$\vthree(T(n,H)) \ge 1$. This half needs only the law, not \S\ref{sec:ladder}.

\emph{At the boundary.} $H$ odd puts $n = 3k+1$, where the explicit exponent is
$0$ and $\vthree(P_k(3k+1)) = 0$ by Theorem~\ref{thm:oddspine}. $H$ even puts
$n = 3k$, where the explicit exponent is $-1$ and $\vthree(P_k(3k)) = 1$ by
Theorem~\ref{thm:evenspine}. Either way $\vthree(T) = 0$.
\end{proof}

\section{The Smith normal form}
\label{sec:snf3}

\begin{theorem}
\label{thm:snf-3power}
For every $N$, the Smith normal form of $T_N := [T(i,j)]_{i,j \le N}$ over $\Z$
is $\mathrm{diag}(3^{e_1},\dots,3^{e_N})$. The cokernel $\Z^N/T_N\Z^N$ is a
finite abelian $3$-group.
\end{theorem}

\begin{proof}
$T(n,H) = 0$ for $H > n$, so $T_N$ is lower triangular and
$\det T_N = \prod_n T(n,n) = \prod_n 3^{\,n-1} = 3^{N(N-1)/2}$, a pure power of
$3$. The invariant factors form a divisibility chain $d_1 \mid \dots \mid d_N$
with $\prod d_i = |\det|$, so any prime $q \ne 3$ dividing some $d_i$ would
divide $d_N$, hence the determinant --- impossible.
\end{proof}

The argument uses only the determinant, so Theorem~\ref{thm:snf-3power} is
unconditional. Counting the nontrivial factors needs Theorem~\ref{thm:activation}, and
with it the congruences of \S\ref{sec:ladder}.

\begin{theorem}
\label{thm:snf}
$T_N$ has exactly $\lceil (N-1)/3\rceil$ nontrivial invariant factors.
\end{theorem}

\begin{proof}
By Theorem~\ref{thm:snf-3power} the invariant factors are $3$-powers, so the
number of nontrivial ones is $N - \mathrm{rank}(T_N \bmod 3) =
\mathrm{nullity}(T_N \bmod 3)$.

By Theorem~\ref{thm:activation}, column $H$ of $T_N$ is identically zero
modulo $3$ if $\lfloor 3H/2\rfloor > N$, and otherwise has its first nonzero
entry in row $\lfloor 3H/2\rfloor$. The map $H \mapsto \lfloor 3H/2\rfloor$ is
strictly increasing, so the nonzero columns have their leading entries in
distinct, increasing rows: they are in echelon position and hence linearly
independent modulo $3$. The null space is therefore spanned exactly by the zero
columns, and
\[
  \mathrm{nullity}(T_N \bmod 3)
  \;=\; \#\{H \le N : \lfloor 3H/2\rfloor > N\}
  \;=\; \left\lceil \frac{N-1}{3} \right\rceil ,
\]
the last step elementary.
\end{proof}

We have verified this directly on the enumerated matrix for every $N \le 36$,
reproducing Table~\ref{tab:snfcount} and the full exponent lists.

\section{Reading along rows}

The results above are statements about columns. Read along rows, the same
triangle has a last nonzero entry in each row, which
Theorem~\ref{thm:activation} already locates for two residue classes of $n$
and not for the third.

\paragraph{The spine is also the last nonzero of its row.} Being to the right of
the spine in row $n$ means $n < \lfloor 3H/2\rfloor$, the same region
Theorem~\ref{thm:activation} proves zero, so the two statements are transposes
of one fact.

\paragraph{Rows $n \equiv 2 \pmod 3$ have no spine cell at all.} For those rows
$\lfloor 3H/2\rfloor$ never equals $n$, so the last nonzero entry is one step
earlier, at the deficit-$2$ cell, whose residue is $2$.

\begin{theorem}[deficit-$2$ residue]
\label{thm:deficit2}
$T(3m+2,\,2m+1) \equiv 2 \pmod 3$ for every $m \ge 0$.
\end{theorem}

\begin{proof}[Method]
Let $D(v) := \sum_m P_{m+1}(3m+2)\,v^m$ be the generating function of the
family, which is a diagonal of $G\,H^n$. Formal Lagrange--B\"urmann over
$\Z/27$, together with the parametrisation $v = u/H^3$, turns the claim
$D = 5 + 18v/(1-v)$ into a rational identity on the mod-$27$ master curve
$E(H,u) = 0$. Since $E$ is monic in $H$, that identity is decided by exact
polynomial division, and the remainder is identically zero.
\end{proof}

So the last-nonzero residue of row $n$ cycles $1, 1, 2$ with $n \bmod 3$: the two
spine theorems supply the first two cases and Theorem~\ref{thm:deficit2} the
third.

\section{Sleeve zeros}
\label{sec:sleeve}

A cell to the left of the spine in row $n$ carries a forced deficit
$d := 3k+1-n$, and $\vthree(P_k(n)) = d$ generically. A \emph{sleeve zero} is a
cell with $\vthree(P_k(n)) > d$. No formula for them is known, and what follows
is a census.

It is complete for the project's data, $n \le 40$. Valuations are read off the
enumerated triangle through $\vthree(P_k(n)) = \vthree(T(n,n-k)) + d$, so no
fitted $P_k$ is needed; fitting would have stopped the census at $k = 17$, while
the proved regime $n \ge 2k+1$ reaches further.

\begin{table}[H]
\centering\small
\begin{tabular}{r l r r}
\toprule
$n$ & sleeve zeros, as $k\!: d \to \vthree(P_k(n))$ & count & longest run \\
\midrule
19 & $7\!: 3\to4$, \; $9\!: 9\to10$ & 2 & 1 \\
37 & $13\!: 3\to5$, \; $16\!: 12\to13$, \; $17\!: 15\to18$ & 3 & 2 \\
38 & $16\!: 11\to13$ & 1 & 1 \\
39 & --- & 0 & 0 \\
40 & $16\!: 9\to10$, \; $18\!: 15\to18$, \; $19\!: 18\to21$ & 3 & 2 \\
\bottomrule
\end{tabular}
\caption{Sleeve zeros in the first row where two occur, and in the last four
rows of the data. The four record events --- first row with two, first adjacent
pair ($n=30$), first three in total ($n=35$), first three in a row ($n=36$) ---
are final for this data: rows 37 and 40 tie the three-in-total record but nothing
reaches four.}
\label{tab:sleeve}
\end{table}

In the final rows the $k = 16$ diagonal spikes at $n = 36, 37$ and $38$, and
again at $n = 40$, skipping only $n = 39$, which is the one sleeve free of zeros
after $n = 31$. The excess $\vthree - d$ there reaches $3$, against a ceiling of
$5$ over all $n$ (at $k=8$, $n=17$).

\subsection{Deficit two and three}
\label{sec:deficit23}

The census above is a census because no unit formula was available past the
spine. Two more deficits now have one, by the same route: fix the tower of
congruences that $H$ and $G$ satisfy modulo a power of $3$, then divide.

At $d = 3$ the statement is
\begin{theorem}
\label{thm:d3}
For every $k \ge 3$,
\[
  P_k(3k-2) \;\equiv\; 27\,r_k \pmod{81},
  \qquad r_k = 2,\,0,\,1 \ \text{ for } k \equiv 0,\,1,\,2 \pmod 3 .
\]
Equivalently, on the deficit-$3$ family $(n,H) = (3k-2, 2k-2)$, where
$T(3k-2,2k-2) = P_k(3k-2)/27$: $T \equiv 2 \pmod 3$ for $k \equiv 0$,
$T \equiv 1 \pmod 3$ for $k \equiv 2$, and $\vthree(P_k(3k-2)) \ge 4$ for
$k \equiv 1$. Since $r_k$ is a unit for $k \not\equiv 1$,
\[
  \vthree\bigl(P_k(3k-2)\bigr) > 3 \iff k \equiv 1 \pmod 3 .
\]
\end{theorem}

The last line is what the sleeve census could only tabulate: at deficit $3$ the
sleeve zeros are exactly the $k \equiv 1 \pmod 3$ cells, for \emph{all} $k$,
rather than for the $k \le 11$ the data reached.

\begin{remark}[the frame this rests on]
Theorem~\ref{thm:d3} is proved on an explicit frame and the frame is not all
theorem. (A1) the diagonal-law shape theorem and the grand form
$P_k(n) = [u^k]G(u)H(u)^n$ --- both proved and Lean-checked. (A2) $H$ congruent
to the fixed point of the mod-$81$ master equation at all orders, derived from
the renewal formalism with completeness from a proved row bound, and
machine-verified against the banked $h_k$ for $k \le 17$. (A3) the corresponding
congruence for $G$, likewise derived and verified to $k \le 17$. (A4) eleven
cluster weights, finite integers from a banked table with independent low-order
enumeration cross-checks. Everything above those four is unconditional algebra:
an inversion lemma, one implicit differentiation, a units inventory, and an exact
polynomial division whose remainder vanishes identically mod $81$. So (A1) and
(A4) are solid, and the weight sits on (A2) and (A3), which are derivations
checked against data rather than proofs.
\end{remark}

Deficit $2$ closes the same way and was done first
(Theorem~\ref{thm:deficit2}); the finer valuations at
$k \equiv 1 \pmod 3$ --- $7, 4, 4, 5$ at $k = 4, 7, 10, 13$ --- are invisible
modulo $81$ and belong to the higher levels of the tower.

\begin{openproblem}[the deficit-$d$ unit formulas, $d \ge 4$]
Theorems~\ref{thm:oddspine}, \ref{thm:evenspine} and~\ref{thm:d3} are the cases
$d = 0, 1, 3$ of a family, with $d = 2$ proved the same way: give the unit part
of $P_k(n)$ at forced deficit $d$ for $d \ge 4$. Each step costs one more level
of the tower --- $d = 3$ needed mod $81$ --- and the open remainder of the
Witt-vector tower of \S\ref{sec:lift} is exactly this.
\end{openproblem}

\section{The 3-adic lift}
\label{sec:lift}

Two further $3$-adic levels of the cubic have been computed.

\begin{itemize}
\item $H^3 - H^2 \equiv 25t \pmod{27}$. The spine cubic lifts two further
  $3$-adic levels, the constant being exactly the pair weight $25$
  ($25 \equiv 7 \bmod 9$ and $\equiv 1 \bmod 3$, which is why the mod-$3$
  statement reads $W^3 = W^2 + t$ with no visible constant).
\item $\bigl(H^3 - H^2 - 25t\bigr)/27 \equiv t\bigl(W(t^3)-1\bigr) \pmod 3$. The
  level-$3$ correction is again $W$, so the same series reappears one level up.
\item Also found: $\bigl(H(t^3) - H^2 - t\bigr)/3 \equiv -(t + t^2 + tW)
  \pmod 3$.
\end{itemize}

These three are measured, not proved.

\section{The $n = 40$ predictions}

The theorems above were stated before the enumeration reached $n = 40$, so the
last rows test them rather than fit them. All three predictions held, cell by
cell, against the final triangle:

\begin{itemize}
\item $T(37,25) \equiv 1 \pmod 3$ (odd spine, $k=12$), and every in-regime cell
  of row $37$ above the spine vanishes modulo $3$. \textbf{Confirmed} --- the
  spine cell is $\equiv 1$, and $T(37,H)$ for $H = 26,\dots,37$ together with
  $T(n,25)$ for $n < 37$ all vanish.
\item $T(39,26) \equiv 1 \pmod 3$ and $P_{13}(39) \equiv 3 \pmod 9$ (even
  spine). \textbf{Confirmed}.
\item Every remaining in-regime cell modulo $3$, via the digit product.
  \textbf{Confirmed}.
\end{itemize}

\section{Related machinery, and why it does not apply}
\label{sec:related}

A reader who works on congruences for combinatorial sequences will reach for two
tools here, and neither reaches this object.

Rowland and Yassawi~\cite{rowlandYassawi2015} compute, for a sequence whose
generating function is the diagonal of a rational power series, a finite
automaton giving the sequence modulo $p^\alpha$ --- which would settle statements
of the kind proved here mechanically. The hypothesis fails: companion paper L4
proves the height generating function is not D-finite, so it is not the diagonal
of a rational function, and no automaton of that provenance exists to compute.
The obstruction is a theorem of ours rather than an absence of effort.

The second tool is the family of methods for the mod-$3^k$ behaviour of
sequences satisfying finite-depth recurrences with polynomial coefficients.
Those apply to a sequence given by such a recurrence; $T(n,H)$ is given by no
recurrence we know, which is the reason this paper works through the master
equation and its fixed point instead.

A literature-priority pass on the spine cubic, the digit product and the Smith
normal form count found no collision (2026-08-18,
\repo{docs/priority-passes-2026-08-18.md}).

\section{Open problems}

\begin{openproblem}[the individual exponents]
Theorem~\ref{thm:snf} counts the nontrivial invariant factors and the
determinant fixes their sum at $N(N-1)/2$, but the individual exponents
$e_i(N)$ are unknown. The mod-$3$ rank gives the count and says nothing about
the higher $3$-adic structure; the lifts of \S\ref{sec:lift} are the natural
attack.
\end{openproblem}

\begin{openproblem}[other primes]
The general spine theorem gives one cubic per prime dividing the drift count,
which for polyhexes is $p = 2$, with the same curve. Does the polyhex triangle
have a mod-$2$ arithmetic as rich as this one, with a digit product, an
activation boundary and a Smith normal form count? Nothing here forbids it and
nobody
has looked.
\end{openproblem}

\section{Reproduction}
\label{sec:repro}

\begin{table}[H]
\centering\small
\begin{tabular}{l l}
\toprule
claim & check \\
\midrule
T1--T5 and the congruences & \repo{experiments/ternary\_spine.py} (15 checks) \\
the 342-cell digit-product mass check & \repo{experiments/ternary\_spine.py} \\
Theorem~\ref{thm:deficit2} & \repo{experiments/deficit2\_proof.py} \\
Table~\ref{tab:sleeve} & \repo{experiments/sleeve\_zero\_census.py} \\
the enumerated triangle & \repo{results/triangle.txt} \\
the underlying law and product form & companion paper L1; \repo{docs/proofs/diagonal-law.md}, \repo{docs/proofs/grand-form.md} \\
the master equation and weights & \repo{results/defect-gas.md} \\
\bottomrule
\end{tabular}
\caption{Where each claim is checked.}
\label{tab:repro}
\end{table}

\bibliographystyle{plain}
\bibliography{shared/refs}

\end{document}
